% ===== PRÉAMBULE CANONIQUE — à coller VERBATIM en tête de CHAQUE .tex =====
\documentclass[11pt,a4paper]{article}
\usepackage[utf8]{inputenc}\usepackage[T1]{fontenc}\usepackage{lmodern}
\usepackage[french]{babel}
\usepackage{amsmath,amssymb,mathtools}\usepackage{icomma}
\usepackage{siunitx}\sisetup{locale=FR}  % \qty{}{}, \si{}, \ang{}
\usepackage{xcolor}\usepackage{colortbl}
\usepackage{tikz}\usepackage{pgfplots}\pgfplotsset{compat=1.18}
\usepackage[siunitx,europeanresistors]{circuitikz} % circuits (élec.)
\usepackage[most]{tcolorbox}\usepackage{enumitem}\usepackage{array,booktabs}
\usepackage[a4paper,margin=2.2cm]{geometry}
\usetikzlibrary{arrows.meta,calc,angles,quotes,positioning,patterns,%
  backgrounds,shapes.geometric,decorations.pathmorphing,decorations.markings,intersections}
\definecolor{primary}{RGB}{222,49,99}\definecolor{primarydark}{RGB}{150,22,55}
\definecolor{defcol}{RGB}{0,114,178}\definecolor{methcol}{RGB}{52,92,148}
\definecolor{excol}{RGB}{0,135,90}\definecolor{attncol}{RGB}{200,40,55}
\tcbset{boxbase/.style={enhanced,breakable,boxrule=0.9pt,arc=2.5pt,
  left=3mm,right=3mm,top=2.3mm,bottom=2.3mm,fonttitle=\bfseries,coltitle=white}}
% --- boîtes TUTORIEL (noms EXACTS, ne pas renommer) ---
\newtcolorbox{cle}[1][]{boxbase,colback=defcol!6,colframe=defcol,
  title={\textsc{À retenir}\ifx\relax#1\relax\relax\else\textmd{ : \itshape #1}\fi}}
\newtcolorbox{methode}[1][]{boxbase,colback=methcol!7,colframe=methcol,sharp corners,
  title={\textsc{Méthode}\ifx\relax#1\relax\relax\else\textmd{ --- \itshape #1}\fi}}
\newtcolorbox{exemplebox}[1][]{boxbase,colback=excol!5,colframe=excol,
  title={\textsc{Exemple}\ifx\relax#1\relax\relax\else\textmd{ : \itshape #1}\fi}}
\newtcolorbox{piege}[1][]{boxbase,colback=attncol!6,colframe=attncol,
  title={\textsc{Piège}\ifx\relax#1\relax\relax\else\textmd{ : \itshape #1}\fi}}
\newtcolorbox{remarque}[1][]{boxbase,colback=black!4,colframe=black!45,
  title={\textsc{Remarque}\ifx\relax#1\relax\relax\else\textmd{ : \itshape #1}\fi}}
% --- boîtes EXERCICE / CORRIGÉ (noms EXACTS, auto-comptées) ---
\newtcolorbox[auto counter]{exo}[1][]{boxbase,colback=primary!4,colframe=primary,
  title={\textsc{Exercice}~\thetcbcounter\ifx\relax#1\relax\relax\else\textmd{ --- \itshape #1}\fi}}
\newtcolorbox[auto counter]{corrige}[1][]{boxbase,colback=excol!4,colframe=excol,
  title={\textsc{Corrigé}~\thetcbcounter\ifx\relax#1\relax\relax\else\textmd{ --- \itshape #1}\fi}}
\newcommand{\vv}[1]{\vec{#1}}\newcommand{\ds}{\displaystyle}
\newcommand{\R}{\mathbb{R}}\newcommand{\N}{\mathbb{N}}\newcommand{\Z}{\mathbb{Z}}
% =========================================================================
\begin{document}
\sisetup{per-mode=symbol}
% ================= DIFFICILE =================

\begin{exo}[la pomme et la Terre --- vrai ou faux ?]
Une pomme de masse $m$ est posée au sol. On s'intéresse aux forces de gravitation
entre la pomme et la Terre. Indique pour chaque proposition si elle est
\textbf{vraie} ou \textbf{fausse}, en justifiant.
\begin{enumerate}[label=\alph*)]
  \item La Terre attire la pomme et la pomme attire la Terre avec deux forces de
        \emph{même} intensité.
  \item Comme la Terre est bien plus massive, elle attire la pomme beaucoup plus fort
        que la pomme n'attire la Terre.
  \item Ces deux forces s'appliquent sur le même corps, donc elles se compensent et
        leur somme est nulle.
  \item La réaction (au sens de la 3\up{e} loi de Newton) de la force gravitationnelle
        exercée par la Terre sur la pomme est la force gravitationnelle exercée par la
        pomme sur la Terre.
\end{enumerate}
\end{exo}

\begin{exo}[deux satellites --- la racine du rapport]
Deux satellites de \emph{même} masse sont en orbite : $A$ a son centre à
$d_A=\qty{8000}{\kilo\metre}$ du centre de la Terre, $B$ à $d_B=\qty{6000}{\kilo\metre}$.
\begin{enumerate}[label=\alph*)]
  \item Établis l'expression du rapport $\dfrac{F_A}{F_B}$ des forces subies.
  \item Calcule la \emph{racine carrée} de ce rapport, $\sqrt{F_A/F_B}$.
  \item Déduis-en lequel des deux satellites subit la force la plus intense.
\end{enumerate}
\end{exo}

\begin{exo}[le Soleil nous attire-t-il plus que la Terre ?]
On compare, sur une personne de masse $m$, la force de gravitation exercée par la
Terre (personne au sol) et celle exercée par le Soleil. Données :
$M_T=\qty{6.0e24}{\kilogram}$, $R_T=\qty{6.4e6}{\metre}$,
$M_S=\qty{2.0e30}{\kilogram}$, distance Terre--Soleil $D=\qty{1.5e11}{\metre}$.
\begin{enumerate}[label=\alph*)]
  \item Exprime le rapport $\dfrac{F_S}{F_T}$ en fonction de $M_S$, $M_T$, $R_T$ et $D$
        (la masse $m$ et $\mathcal{G}$ doivent se simplifier).
  \item Calcule sa valeur numérique.
  \item Qui, du Soleil ou de la Terre, attire le plus fortement la personne ?
\end{enumerate}
\end{exo}

\begin{exo}[combien de propositions correctes ?]
Deux corps s'attirent avec une force $F$. Parmi les affirmations suivantes,
indique celles qui sont \textbf{correctes}.
\begin{enumerate}[label=\alph*)]
  \item Si l'on double \emph{chacune} des deux masses, la force est multipliée par $4$.
  \item Si l'on triple la distance, la force est divisée par $3$.
  \item Si l'on double une masse et que l'on double aussi la distance, la force est
        divisée par $2$.
  \item Si l'on divise la distance par $3$, la force est multipliée par $9$.
\end{enumerate}
Combien de ces propositions sont correctes ?
\end{exo}

\end{document}
